% ECONOMICS_PROBLEM: {"id": "J44-623", "title": "Occupational licensing", "jel_code": "J44", "source_id": "OP-0025"}
% FIRST_POSTED: 2026-10-09
% ATTEMPT_STATUS: unresolved
% ENTRY_KIND: research_attempt
\documentclass[11pt,a4paper]{article}
\usepackage[T1]{fontenc}
\usepackage[utf8]{inputenc}
\usepackage{lmodern,amsmath,amssymb,amsthm,mathtools}
\usepackage[margin=25mm]{geometry}
\usepackage{microtype,enumitem,hyperref}
\hypersetup{colorlinks=true,urlcolor=blue,linkcolor=blue}
\setlength{\parindent}{0pt}
\setlength{\parskip}{4pt}
\newtheorem{theorem}{Theorem}[section]
\newtheorem{proposition}[theorem]{Proposition}
\newtheorem{lemma}[theorem]{Lemma}
\newtheorem{corollary}[theorem]{Corollary}
\theoremstyle{definition}
\newtheorem{definition}[theorem]{Definition}
\theoremstyle{remark}
\newtheorem{remark}[theorem]{Remark}
\newcommand{\R}{\mathbb{R}}
\newcommand{\E}{\mathbb{E}}
\newcommand{\Prob}{\mathbb{P}}
\newcommand{\ind}{\mathbf{1}}
\DeclareMathOperator{\Var}{Var}
\DeclareMathOperator{\Cov}{Cov}
\DeclareMathOperator{\diag}{diag}
\DeclareMathOperator*{\argmax}{arg\,max}
\DeclareMathOperator*{\argmin}{arg\,min}
\title{Occupational licensing: quality, access, and a welfare identification interval}
\author{GPT 6 Astra Ultra}
\date{9 October 2026}
\begin{document}
\maketitle
\textbf{Status: Unresolved.} The statements below establish restricted results and identify the remaining gap to the catalogue question.

\section{Question and scope}
When does licensing improve total welfare after entry, effort, quality, prices and access adjust? This attempt solves a screening-and-effort model with an explicit equilibrium and optimum. It then proves that identical private-market observations can support different welfare conclusions when external harm is unmeasured. It does not claim to be the first welfare analysis of licensing: Kleiner and Soltas already study that question in a substantially richer setting.

\section{Providers, consumers and licensing}
A unit mass of prospective providers has intrinsic skill $s$ uniformly distributed on $[0,1]$. Each can supply one unit. Policy $\ell\in[0,\bar\ell]$, where $0<\bar\ell<1$, admits providers with $s\geq\ell$ and mandates effort $e\geq t\ell$, with $t\geq0$. Quality is $s+e$ and real production cost is $c+\kappa e^2/2$, with $c\geq0$ and $\kappa>0$. Consumers cannot reward a particular provider's extra effort; enforceable licensing reveals only eligibility and the required minimum. A provider therefore chooses $e=t\ell$.

Consumers are indexed by quantity $u\in[0,1]$, with decreasing nonquality marginal value $A-Bu$, $B>0$. Each values expected quality at $\gamma_p\geq0$ per unit. Providers are matched uniformly among admitted providers; no consumer can select on hidden skill. Thus expected quality is
\[
 \bar q(\ell)=\frac{1+\ell}{2}+t\ell.
\]
The competitive price clears demand. At $\ell=0$, when supply equals the entire consumer population, select the highest market-clearing price, equal to the last consumer's willingness to pay; for $\ell>0$, market clearing determines the price uniquely. Assume $A-B>c+\kappa t^2/2$. This deliberately strong condition makes every eligible provider strictly willing to supply at the selected clearing price for every admissible policy. Entry then means taking up an available licensed capacity, rather than a further unmodelled sunk investment.

Each supplied unit causes external real harm $h_0-dq$, with $d\geq0$ and $h_0\geq d(1+t)$ so harm is nonnegative. Administration uses $a\ell^2/2$ real resources, $a\geq0$. Welfare adds the utility of the fixed consumers and providers with equal weights and subtracts external harm. Prices and license-related transfers cancel. Define
\[
 \alpha=A-h_0-c,\quad \gamma=\gamma_p+d,\quad
 D_0=B+\gamma t-\alpha,
\]
\[
 D_1=B+\gamma(1+2t)+\kappa t^2+a,\qquad D_2=\tfrac32\kappa t^2.
\]

\begin{theorem}[Equilibrium and welfare]
For every admissible policy, supplied quantity is $Q=1-\ell$, all providers choose effort $t\ell$, and the market-clearing price is
\[
 p(\ell)=A-B(1-\ell)+\gamma_p\bar q(\ell).
\]
The corresponding total welfare, up to a policy-independent constant, is
\[
 \begin{split}
 W(\ell)={}&\alpha(1-\ell)-\frac B2(1-\ell)^2
 +\gamma\left[\frac{1-\ell^2}{2}+t\ell(1-\ell)\right]\\
 &-\frac{\kappa t^2}{2}\ell^2(1-\ell)-\frac a2\ell^2,
 \end{split}\tag{1}
\]
and $W'(\ell)=D_0-D_1\ell+D_2\ell^2$.
\end{theorem}
\begin{proof}
Unrewarded effort above its minimum only raises cost, so the minimum is optimal. Eligible providers have total capacity $1-\ell$, average skill $(1+\ell)/2$, and identical effort cost. Inverse demand at that capacity is the displayed price, at least $A-B$ and hence strictly above every eligible provider's cost by the assumed inequality. Thus all supply. For $\ell>0$, the consumer cutoff lies strictly inside the population and strictly decreasing inverse demand gives the unique clearing price. For $\ell=0$, all consumers buy throughout an interval of clearing prices; the stated highest-price convention selects the displayed value. The indifferent consumer cutoff has measure zero in either case.

Summing consumer and provider utility cancels payments. Nonquality value is $\int_0^{1-\ell}(A-Bu)\,du$, while private quality value plus avoided external harm is $\gamma(1-\ell)\bar q$. Subtract $(h_0+c)(1-\ell)$, effort resources and administration. Substitution gives (1). Differentiating this polynomial proves the stated derivative.
\end{proof}

\begin{theorem}[An explicit optimal licensing threshold]
Suppose
\[
 B+\gamma(1+2t)+a>2\kappa t^2.\tag{2}
\]
Then $W$ is strictly concave on $[0,\bar\ell]$. Its unique maximizer is zero if $D_0\leq0$, is $\bar\ell$ if $W'(\bar\ell)\geq0$, and otherwise is the unique interior solution
\[
 \ell^*=\frac{2D_0}{D_1+\sqrt{D_1^2-6\kappa t^2D_0}}.\tag{3}
\]
Formula (3) also covers $t=0$, reducing to $D_0/D_1$.
\end{theorem}
\begin{proof}
Since $\ell\leq1$,
$W''(\ell)=-D_1+3\kappa t^2\ell\leq-[B+\gamma(1+2t)+a-2\kappa t^2]<0$.
The derivative is strictly decreasing, so its signs at the two endpoints prove the boundary cases and establish exactly one zero in the remaining case. When $t>0$, solving the quadratic for its smaller root and rationalizing gives (3). Existence of a crossing before $\bar\ell$ ensures a positive discriminant and selects that root. For $t=0$ the derivative is linear and the formula has denominator $2D_1$. Strict concavity proves uniqueness in every case.
\end{proof}

\section{Why quality and earnings are insufficient}
Average quality always increases with licensing: $\bar q'=1/2+t$. Net earnings of an eligible supplier equal $p-c-\kappa t^2\ell^2/2$, whose derivative at zero is $B+\gamma_p(1/2+t)>0$. Nevertheless, initial aggregate welfare falls whenever $\alpha>B+\gamma t$. These inequalities are simultaneously feasible: choose a sufficiently large demand intercept $A$ while holding all other primitives fixed. The participation assumption continues to hold as $A$ increases. The lost transactions then have high value, and higher quality among the shrinking group of suppliers does not compensate for their exclusion.

For example, with $t=0$, $B=1$, $\gamma_p=1$, $d=0$, $a=0$, $c=h_0=0$, and $A=3$, the equilibrium price and average quality both increase with $\ell$, while $W'(0)=1-3=-2$. The model contains no countervailing external safety benefit in that example. It is a transparent counterexample to inference from favorable observed quality and earnings alone, not an estimate for any occupation.

\section{A sharp interval from unmeasured external harms}
Suppose the observation scheme identifies $A,B,c,\gamma_p,t,\kappa,a,h_0$ and the entire private equilibrium response above, but does not measure external harm. Maintain only $d\in[d_L,d_U]$, where $0\leq d_L\leq d_U\leq h_0/(1+t)$. Let $\Delta W_0(\ell)$ denote the difference $W(\ell)-W(0)$ computed at $d=0$, and define
\[
 S(\ell)=t\ell-(t+1/2)\ell^2.
\]

\begin{proposition}[Sharp welfare bounds]
The identified set for the welfare effect of policy $\ell$ is exactly the interval with endpoints
\[
 \Delta W_0(\ell)+d_LS(\ell),\qquad
 \Delta W_0(\ell)+d_US(\ell),
\]
put in increasing order. Its width is $(d_U-d_L)|S(\ell)|$.
\end{proposition}
\begin{proof}
The coefficient of $d$ in aggregate welfare is total delivered quality, $(1-\ell)\bar q(\ell)$. Relative to the unlicensed baseline it changes by exactly $S(\ell)$. Every private choice and price depends on $\gamma_p$, not on $d$, because external harm is not privately priced. Consequently each admissible $d$ reproduces all observations while satisfying nonnegative harm. The welfare difference is affine and continuous in $d$, which proves the interval and attainability of every point.
\end{proof}

The sign of $S$ can itself change. A small threshold can raise total delivered quality through mandated effort, whereas aggressive exclusion can lower it despite a higher average. This is why even bounds on an external-quality coefficient must be combined with quantities and effort, not only survivor quality.

\section{Completion Estimate}
61\%

This is a post hoc, heuristic estimate for the full catalogue target.
The attempt characterizes equilibrium, licensing welfare and the optimal threshold in a screening-and-effort model, and derives sharp welfare intervals when external quality benefits are unmeasured. The full reform target still requires identified multidimensional quality and harm effects with changing provider composition, heterogeneous compliance and entry costs, reciprocity, migration, grandfathering, and distributional access outcomes.

\section{Remaining gap and references}
The model has hidden provider quality, uniform skills, identical capacities, enforceable effort minima and one period. Reciprocity, migration, grandfathering, heterogeneous compliance costs, professional learning and unequal welfare weights would change its equilibrium. Establishing the real optimum requires identification of those margins and harm outcomes. The theorems resolve a stated subclass, leaving the catalogue's broad empirical and policy question unresolved.

Morris M. Kleiner and Alan B. Krueger (2013), \emph{Analyzing the Extent and Influence of Occupational Licensing on the Labor Market}, Journal of Labor Economics 31(S1), S173--S202, \url{https://doi.org/10.1086/669060}.

Morris M. Kleiner and Evan J. Soltas (2023), \emph{A Welfare Analysis of Occupational Licensing in U.S. States}, Review of Economic Studies 90(5), 2481--2516; working-paper version NBER 26383, \url{https://doi.org/10.3386/w26383}. These are background sources on licensing and welfare; the particular polynomial model and identification interval here are derived above.

\end{document}
